%! Characteristics of Exponential Functions — Unit 7, Lesson 7.0 — Homework (Answer Key)
\documentclass[10pt]{article}
\usepackage{algebra2-article}
\usepackage{algebra2-key}   % pulls in -boxes; do NOT also load -boxes
\newcommand{\TallMath}[1]{$\displaystyle #1\rule[-1.4em]{0pt}{3.2em}$}
\newcommand{\lgrid}[4]{%
  \draw[step=1, linegray!40, very thin] (#1,#3) grid (#2,#4);
  \draw[->, semithick, charcoal] (#1,0)--(#2,0) node[below right, font=\scriptsize]{$x$};
  \draw[->, semithick, charcoal] (0,#3)--(0,#4) node[left, font=\scriptsize]{$y$};}
% #2 is ln(b):  ln2 = 0.693147   ln3 = 1.098612
\newcommand{\expcurve}[4]{\draw[very thick, forest, domain=#3:#4, samples=120, smooth]
  plot (\x,{#1*exp((#2)*(\x))});}

\begin{document}

\pageheader{Unit 7, Lesson 7.0}{Homework --- Answer Key}
\namedateperiod

\vspace{0.10in}

\begin{notesbox}{Practice}
Show how you read each feature. Write every domain and range in \textbf{interval notation}.

\begin{enumerate}[leftmargin=1.9em, itemsep=12pt]
  \item \textbf{From the equation only.} Use $a$ and $b$ --- no graphing.
    \begin{center}
    \renewcommand{\arraystretch}{1.6}
    \begin{tabularx}{\linewidth}{@{}l X X X X@{}}
    \hline
    \textbf{Function} & \textbf{Base $b$} & \textbf{Growth or decay?} & \textbf{$y$-intercept} & \textbf{Asymptote approached on the left or right?} \\
    \hline
    $y=6\cdot 2^{x}$ & \ans{$2$} & \ans{growth} & \ans{$(0,6)$} & \ans{left} \\
    $y=6\left(\frac12\right)^{x}$ & \ans{$\frac12$} & \ans{decay} & \ans{$(0,6)$} & \ans{right} \\
    $y=\left(\frac45\right)^{x}$ & \ans{$\frac45$} & \ans{decay} & \ans{$(0,1)$} & \ans{right} \\
    $y=0.4(3)^{x}$ & \ans{$3$} & \ans{growth} & \ans{$(0,0.4)$} & \ans{left} \\
    \hline
    \end{tabularx}
    \end{center}
    All four have the same domain, the same range, and the same horizontal asymptote. Write them
    once: domain \ans{$(-\infty,\infty)$}, range \ans{$(0,\infty)$}, asymptote \ans{$y=0$}.

  \item \textbf{A growth graph.} The graph is $m(x) = 4\cdot 2^{x}$.
    \begin{center}
    \begin{tikzpicture}[scale=0.42]
      \lgrid{-4}{3}{-1}{10}
      \begin{scope}
        \clip (-4,-1) rectangle (3,10);
        \expcurve{4}{0.693147}{-4}{1.32}
      \end{scope}
      \draw[navy, dashed, semithick] (-4,0)--(3,0);
      \fill[charcoal] (-2,1) circle (3pt) node[below right, font=\scriptsize]{$(-2,1)$};
      \fill[charcoal] (0,4) circle (3pt) node[left, font=\scriptsize]{$(0,4)$};
      \fill[charcoal] (1,8) circle (3pt) node[left, font=\scriptsize]{$(1,8)$};
    \end{tikzpicture}
    \end{center}
    $a=$ \ans{4} \quad $b=$ \ans{2} \quad Growth or decay? \ans{growth} \\[3pt]
    Domain: \ans{$(-\infty,\infty)$} \quad Range: \ans{$(0,\infty)$} \quad Horizontal asymptote: \ans{$y=0$} \\[3pt]
    $y$-intercept: \ans{$(0,4)$} \quad $x$-intercept: \ans{none} \quad
    Increasing or decreasing? \ans{increasing} \\[3pt]
    Absolute extrema: \ans{none} \\[3pt]
    End behavior: as $x\to+\infty$, $m(x)\to$ \ans{$+\infty$}; \; as $x\to-\infty$,
    $m(x)\to$ \ans{$0$}

  \item \textbf{A decay graph.} The graph is $n(x) = 9\left(\frac13\right)^{x}$.
    \begin{center}
    \begin{tikzpicture}[scale=0.42]
      \lgrid{-1}{6}{-1}{10}
      \begin{scope}
        \clip (-1,-1) rectangle (6,10);
        \expcurve{9}{-1.098612}{-0.09}{6}
      \end{scope}
      \draw[navy, dashed, semithick] (-1,0)--(6,0);
      \fill[charcoal] (0,9) circle (3pt) node[right, font=\scriptsize]{$(0,9)$};
      \fill[charcoal] (1,3) circle (3pt) node[right, font=\scriptsize]{$(1,3)$};
      \fill[charcoal] (2,1) circle (3pt) node[above right, font=\scriptsize]{$(2,1)$};
    \end{tikzpicture}
    \end{center}
    $a=$ \ans{9} \quad $b=$ \ans{$\frac13$} \quad Growth or decay? \ans{decay} \\[3pt]
    Domain: \ans{$(-\infty,\infty)$} \quad Range: \ans{$(0,\infty)$} \quad Horizontal asymptote: \ans{$y=0$} \\[3pt]
    $y$-intercept: \ans{$(0,9)$} \quad $x$-intercept: \ans{none} \quad
    Increasing or decreasing? \ans{decreasing} \\[3pt]
    End behavior: as $x\to+\infty$, $n(x)\to$ \ans{$0$}; \; as $x\to-\infty$,
    $n(x)\to$ \ans{$+\infty$} \\[3pt]
    Confirm $b$ from the labeled points: $\frac{n(1)}{n(0)}=\frac{3}{9}=$ \ans{$\frac13$} \\[3pt]
    Problems 2 and 3 have the \emph{same} horizontal asymptote. Which side does each graph approach
    it on, and what in the equation decides that? \ans{$m$ approaches $y=0$ on the \emph{left}
    and $n$ approaches it on the \emph{right}; the base decides --- $b>1$ shrinks toward the
    asymptote as $x$ goes to $-\infty$, while $0<b<1$ shrinks toward it as $x$ goes to $+\infty$.}

  \item \textbf{Name the family from the table.} Each table has evenly spaced $x$-values. Label each
    one \emph{linear}, \emph{quadratic}, or \emph{exponential}, and give the constant that proves it.
    \begin{center}
    \renewcommand{\arraystretch}{1.25}
    \begin{tabular}{|c|c|c|c|c|c|}
    \hline
    $x$ & $0$ & $1$ & $2$ & $3$ & $4$ \\ \hline
    A: $y$ & $5$ & $8$ & $11$ & $14$ & $17$ \\ \hline
    B: $y$ & $5$ & $10$ & $20$ & $40$ & $80$ \\ \hline
    C: $y$ & $3$ & $4$ & $7$ & $12$ & $19$ \\ \hline
    \end{tabular}
    \end{center}
    A is \ans{linear} because \ans{the difference is a constant $+3$} \\[3pt]
    B is \ans{exponential} because \ans{the ratio is a constant $2$} \\[3pt]
    C is \ans{quadratic} because \ans{the differences $1,3,5,7$ are not constant, but the second differences are a constant $+2$}

  \item \textbf{Explain.} In Unit 6, the endpoint of $y=\sqrt{x}$ gave the function an absolute
        minimum of $0$. An exponential graph like $y=2^{x}$ also has a lowest edge at $y=0$ --- but
        it has \emph{no} absolute minimum. Explain the difference, using the words
        \textbf{endpoint}, \textbf{asymptote}, and \textbf{reaches}.
        \ansline{An \textbf{endpoint} is a point \emph{on} the graph: $\sqrt{0}=0$, so $y=\sqrt{x}$
        actually \textbf{reaches} the value $0$ at $x=0$, and a value the function reaches and never
        goes below is an absolute minimum.\\[3pt]
        An \textbf{asymptote} is a line the graph only approaches. There is no input that makes
        $2^{x}$ equal $0$ --- the outputs get smaller and smaller but stay positive --- so $0$ is
        never \textbf{reached}. It is the boundary of the range, not a value of the function, which
        is why the range is $(0,\infty)$ and there is no minimum.}
\end{enumerate}
\end{notesbox}

\begin{extensionbox}
\textbf{Reading a model in context.} A new laptop costs \$$1200$. Like most electronics, it loses
about $25\%$ of its value every year, so its value after $t$ years is
\[
  V(t) = 1200(0.75)^{t} \ \text{ dollars.}
\]
\begin{center}
\renewcommand{\arraystretch}{1.25}
\begin{tabular}{|c|c|c|c|c|c|}
\hline
$t$ (years) & $0$ & $1$ & $2$ & $3$ & $4$ \\ \hline
$V(t)$ (dollars) & $1200$ & $900$ & $675$ & $506.25$ & $379.69$ \\ \hline
\end{tabular}
\end{center}
\begin{enumerate}[label=(\alph*), leftmargin=1.8em, itemsep=5pt]
  \item Identify $a$ and $b$, say whether this is growth or decay, and explain what each number means
        about the laptop. Why is $b=0.75$ rather than $0.25$?
  \item Confirm the constant ratio using two consecutive values from the table.
  \item Find $V(2)$ and write one sentence interpreting it in context.
  \item A friend says ``it loses \$$300$ a year.'' Use the table to show why that is wrong, and name
        what is actually constant.
  \item The model's horizontal asymptote is $V=0$, and there is no $t$-intercept. State what that
        claims about the laptop, and say whether you believe it.
  \item What is the domain of $V$ algebraically? What is a sensible domain \emph{in context}?
\end{enumerate}
\ansline{(a) $a=1200$, $b=0.75$; \textbf{decay}, since $0<b<1$. $a$ is the price when the laptop is
new ($t=0$); $b=0.75$ says $75\%$ of the previous year's value survives. It is $0.75$, not $0.25$,
because $b$ is the fraction that \emph{remains}, not the fraction lost --- losing $25\%$ leaves
$100\%-25\%=75\%$.\\[3pt]
(b) $\frac{900}{1200}=0.75$ and $\frac{675}{900}=0.75$ --- the same ratio, which is $b$.\\[3pt]
(c) $V(2)=675$: two years after purchase, the laptop is worth about \$$675$.\\[3pt]
(d) The drops are $\$300$, then $\$225$, then $\$168.75$, then $\$126.56$ --- not constant. What is
constant is the \emph{ratio}: it loses $25\%$ of whatever it is currently worth, and $25\%$ of a
smaller value is a smaller dollar amount.\\[3pt]
(e) It claims the laptop's value shrinks forever without ever hitting \$$0$ --- it is always worth
\emph{something}, no matter how old. That is not realistic: at some point the resale value is
genuinely \$$0$, and long before that the predicted value is less than a penny. The model is useful
for the first several years and meaningless in the far tail.\\[3pt]
(f) Algebraically the domain is all real numbers, since any real exponent is legal. In context, $t\ge
0$ --- negative $t$ would mean before the laptop existed --- so a sensible domain is $[0,\infty)$,
and realistically only the first several years.}
\end{extensionbox}

\begin{spiralbox}
\textbf{Coming next --- Lesson 7.1: Exponential Growth \& Decay.} Today you read $a$ and $b$ off a
graph that was handed to you. Next you will \emph{build} $y=ab^{x}$ yourself --- from a table, from
two points, and from a sentence like ``the population grows $8\%$ a year'' or ``the car loses $15\%$
of its value.'' The whole lesson turns on one translation: a percent is not a base. ``Up $8\%$''
means $b=1.08$, not $b=0.08$ --- and ``down $15\%$'' means $b=0.85$, not $b=1.15$. You already met
the idea today: the medication kept $70\%$ each hour, so $b$ was $0.7$.
\end{spiralbox}

\begin{teachernote}
\textbf{Problem 1, last row} ($y=0.4(3)^{x}$) is the one to check: $a$ is a decimal less than $1$ and
the base is greater than $1$, so a student running on ``small number $\Rightarrow$ decay'' will
answer wrong. Growth versus decay is read off $b$ alone, never off $a$.

\textbf{Problem 3's last part} and \textbf{Problem 5} carry the lesson. Problem 3 checks that
students can attach the asymptote to a \emph{side} rather than reciting ``$y=0$''; Problem 5 is the
Unit 6 $\rightarrow$ Unit 7 hinge and is worth grading closely. A complete answer must say that
$\sqrt{0}=0$ is an actual output while no input makes $2^{x}=0$; ``one is a point and one is a line''
alone is not enough.

\textbf{Problem 4} previews Lesson 7.5, where the same three-way test is applied to real scatterplot
data. Table C is the one students misname --- rapid-looking growth reads as ``exponential'' until
they check the ratios ($\frac43$, $\frac74$, $\frac{12}{7}$ --- not constant).

\textbf{Extension (d) and (e)} are the same two beats as the Activity's Tier E drug model, on purpose.
If students did Tier E in class, this is consolidation and can be graded quickly; if they did not,
this is where the constant-ratio idea meets context for the first time and deserves a full read.
Part (a)'s ``why $0.75$ and not $0.25$'' is the direct set-up for Lesson 7.1's percent-to-factor
translation.
\end{teachernote}

\end{document}
